% ============================================================
% main.tex — the test itself
%
% This file is structure and content. All typographic
% choices live in examstyle.sty, which also loads babel.
%
% Compile with pdflatex:
%   latexmk -pdf main
% (two runs: "Page 2 of 5" needs the page count from the
% previous pass — latexmk handles it by itself.)
% ============================================================


\documentclass[12pt]{article}


% ============================================================
% LOAD THE EXAM STYLE
%
% [coverpage]: title block and instructions on a page of
% their own without a page number; Part 1 begins on page 1.
% For full-day tests and tests with two parts. Without the
% option the title block sits at the top of page 1 with the
% exercises straight below — for smaller tests.
%
% [booklet]: like [coverpage], but laid out for duplex
% printing and folding — blank back on the cover, every part
% on a right-hand page, and a page count divisible by 4.
%
% [statuscheck]: title, subtitle and date — no fields. Split
% the exercises with \level[...]{1}, {2}, {3} instead of
% \exampart.
%
% [indent]: paragraph indentation instead of space between
% paragraphs. The options can be combined.
% ============================================================
\usepackage[coverpage]{examstyle}
% \usepackage[booklet]{examstyle}
% \usepackage{examstyle}
% \usepackage[statuscheck]{examstyle}
% \usepackage[indent]{examstyle}


% ============================================================
% CUSTOM PACKAGES
%
% Add project-specific packages here.
% Do not put packages anywhere else in the file.
% ============================================================

% --- Mathematics (standard) ---
\usepackage{amsmath, amssymb}

% --- Vectors: \vv{F}, flat arrow over the symbol. Indices outside: \vv{F}_1 ---
\usepackage[e]{esvect}

% --- Units: \qty{9.81}{\metre\per\second\squared} ---
\usepackage{siunitx}
% The two house choices (m/s^2, combined uncertainty) live in
% the style file. Override here if needed:
% \sisetup{per-mode=power}

% --- Figures ---
\usepackage{graphicx}
\graphicspath{{figures/}}

% --- Other ---
\usepackage{tikz}


% ============================================================
% CUSTOM COMMANDS
% ============================================================

% Upright d:
% \newcommand{\dd}{\mathrm{d}}
% \newcommand{\dv}[2]{\frac{\dd #1}{\dd #2}}


% ============================================================
% METADATA
%
% Empty fields are shown as empty (not omitted). \namelabel
% swaps "Name" for e.g. "Candidate number".
% ============================================================
\title{Physics 1 Test}
\subtitle{Chapters 3 and 4: Forces and motion}
\date{Thursday 15 October 2026}
\duration{90 minutes}
\aids{Part 1: none. Part 2: all, except communication}
% \namelabel{Candidate number}


% ============================================================
% DOCUMENT
% ============================================================
\begin{document}

\maketitle

% With [coverpage] everything from here to the first \exampart stays on the cover.
Answer all the exercises. Show your working and justify your
answers; a correct answer without justification does not earn
full credit. Use figures where helpful.

Part 1 is handed in before you begin Part 2. The test has
5 exercises and is worth 20 points in total.


% ============================================================
% PART 1
%
% With [coverpage] the first \exampart closes the cover and
% begins on page 1. Drop the \exampart lines entirely on a
% smaller test without parts (and without [coverpage]) — the
% exercises count the same either way.
% ============================================================
\exampart[Without aids]{1}


\begin{exercise}[Forces on an incline \points{4}]
A block of mass $m$ rests on an incline making an angle
$\theta$ with the horizontal.

\begin{subproblems}
  \item Draw a figure and mark all the forces acting on the block.
  \item Show that the normal force is $N = mg\cos\theta$.
  \item How large must the friction force be, at minimum, for
        the block to stay at rest?
\end{subproblems}
\end{exercise}


\begin{exercise}[Equations of motion \points{4}]
A ball is thrown straight up with speed $v_0 = \qty{12}{\metre\per\second}$.
Neglect air resistance, and use $g = \qty{9.8}{\metre\per\second\squared}$.

\begin{subproblems}
  \item How high does the ball go?
  \item How long does it take before the ball is back at its starting point?
\end{subproblems}
\end{exercise}


\begin{exercise}[Vectors \points{2}]
Two forces $\vv{F}_1$ and $\vv{F}_2$ act at the same point.
$\vv{F}_1$ has magnitude \qty{3.0}{\newton} and points east;
$\vv{F}_2$ has magnitude \qty{4.0}{\newton} and points north.
Find the magnitude and direction of the resultant force.
\end{exercise}


% ============================================================
% PART 2
%
% \clearpage if Part 1 is to be handed in before Part 2
% begins. To restart the exercises at 1 here:
% \setcounter{exerciseinner}{0} after \exampart.
% ============================================================
\clearpage
\exampart[With aids]{2}


\begin{exercise}[The braking car \points{6}]
A car is travelling at \qty{80}{\kilo\metre\per\hour} when the
driver notices an obstacle \qty{50}{\metre} ahead. The reaction
time is \qty{0.8}{\second}, and the brakes give a constant
deceleration of \qty{6.5}{\metre\per\second\squared}.

\begin{subproblems}
  \item How far does the car travel during the reaction time?
  \item Does the car manage to stop before the obstacle? Show your working.
\end{subproblems}

The figure below shows a velocity--time graph of the motion.

\begin{figure}[ht]
  \centering
  \begin{tikzpicture}[>=stealth, x=1.1cm, y=0.9cm]
    \draw[->,thick] (0,0) -- (5.2,0) node[right,font=\small] {$t$ (s)};
    \draw[->,thick] (0,0) -- (0,3.4) node[above,font=\small] {$v$ (m/s)};
    \draw[gray,thin,dashed] (0.8,0) -- (0.8,2.5);
    \node[below,font=\small,gray] at (0.8,0) {$0.8$};
    \node[below,font=\small,gray] at (4.2,0) {$t_{\text{stop}}$};
    \node[left,font=\small] at (0,2.5) {$22.2$};
    \fill[darkorange!15] (0,0) -- (0,2.5) -- (0.8,2.5) -- (4.2,0) -- cycle;
    \draw[darkolive,thick] (0,2.5) -- (0.8,2.5) -- (4.2,0);
  \end{tikzpicture}
  \caption{Velocity as a function of time. The area under the
           graph is the distance.}
  \label{fig:vt}
\end{figure}

\begin{subproblems}[resume]
  \item Use figure~\ref{fig:vt} to explain how the stopping
        distance can be read off the graph.
\end{subproblems}
\end{exercise}


\begin{exercise}[Numerical solution \points{4}]
The program below simulates the fall of a ball with air
resistance using Euler's method.

\begin{python}
import numpy as np

m, g, k = 0.10, 9.81, 0.02     # mass, g, drag
v, t, dt = 0.0, 0.0, 0.01

while t < 5.0:
    a = g - (k / m) * v**2      # acceleration (m/s^2)
    v = v + a * dt
    t = t + dt

print(f"Speed after {t:.1f} s: {v:.2f} m/s")
\end{python}

\begin{subproblems}
  \item Explain what line 7 does, and which physical law it is based on.
  \item What will the program print if $k$ is set to zero?
        Justify your answer without running the program.
  \item Change the program so that it also computes how far
        the ball has fallen.
\end{subproblems}
\end{exercise}

\separator

\begin{center}
\sffamily Good luck!
\end{center}

\end{document}
